% ============================================================
%  高中数学题库 LaTeX 模板（MathCyclus格式）
%  适用于各种考试题库的排版
%  说明：
%    - 采用 12pt 字号 + 2.5cm 边距
%    - 使用 ctexart 支持中文
%    - 支持选择题、填空题、解答题
%    - 包含答案和解析
%    - 符合 MathCyclus 题库管理系统格式
% ============================================================

\documentclass[12pt, a4paper, oneside]{ctexart}

% ── 1. 数学与链接 ──
\usepackage{amsmath, amsthm, amssymb}
\usepackage[bookmarks=true, colorlinks, citecolor=blue, linkcolor=black]{hyperref}

% ── 2. 字体方案 ──
\usepackage{newtxmath}

% ── 3. 页面布局 ──
\usepackage[a4paper, margin=2.5cm, footskip=1cm]{geometry}

% ── 4. 页眉页脚 ──
\usepackage{fancyhdr}
\usepackage{lastpage}
\pagestyle{fancy}
\fancyhf{}
\fancyfoot[C]{题库第\thepage 页（共\pageref{LastPage}页）}
\renewcommand{\headrulewidth}{0pt}
\renewcommand{\footrulewidth}{0pt}

% ── 5. 表格与插图 ──
\usepackage{graphicx}
\usepackage{adjustbox}
\usepackage{diagbox}
\usepackage{makecell}
\usepackage{caption}
\usepackage{float}
\usepackage{tikz}

% ── 6. 选择题选项（tasks 环境） ──
\usepackage{tasks}
\settasks{
    label       = \Alph*.,
    label-width = 1.8em,
    item-indent = 2.4em,
    label-offset = 0.5em,
    column-sep  = 2em,
    before-skip = 0pt,
    after-skip  = 0pt
}

% ── 7. 大题编号（enumitem 体系） ──
\usepackage[shortlabels]{enumitem}
\newlist{examenum}{enumerate}{3}
\setlist[examenum,1]{label=\arabic*., leftmargin=2em, itemsep=0.3em, parsep=0em}
\setlist[examenum,2]{label=(\arabic*), leftmargin=1.5em, itemsep=0.1em, parsep=0em}
\setlist[examenum,3]{label=(\roman*), leftmargin=1.5em, itemsep=0.1em, parsep=0em}

% ── 8. 自定义命令 ──

% 圈号数字
\newcommand{\mycircled}[1]{%
  \tikz[baseline=(char.base),outer sep=0pt]{%
    \node[draw,circle,inner sep=0.5pt,minimum size=1.4em,
          line width=0.4pt,font=\zihao{-5}] (char) {#1};%
  }%
}

% 旋转平行符号
\newcommand{\Parallel}{\raisebox{0.1ex}{\rotatebox[origin=c]{-20}{$\parallel$}}}

% 下划线填空
\newcommand{\blank}{\underline{\hspace{2cm}}}

% 答案命令
\newcommand{\answer}[1]{\textbf{答案：}#1}

% 解析命令
\newcommand{\analysis}[1]{\textbf{解析：}#1}

% 详解命令
\newcommand{\explanation}[1]{\textbf{详解：}#1}

% ── 9. 大标题辅助命令 ──
\newcommand{\subjecttitle}[1]{{\fontsize{24pt}{22pt}\selectfont\centering\textbf{#1}\par}}
\newcommand{\subtitle}[1]{{\fontsize{16pt}{16pt}\selectfont\centering #1\par}}

% ============================================================
%  正文开始
% ============================================================
\begin{document}

% ── 题库标题区 ──
\begin{center}
    \LARGE{\textbf{考试题库}}
\end{center}

\vspace{0.5em}

% ── 注意事项 ──
\noindent\textbf{注意事项}：

1．本题库包含选择题、填空题、解答题三种题型。

2．选择题和填空题答案直接标注在题目后，解答题提供详细解析。

3．使用本题库时，请先独立完成题目，再对照答案和解析。

\vspace{1em}

% ══════════════════════════════════════════
%  选择题部分
% ══════════════════════════════════════════

\noindent\textbf{选择题：本题共8小题，每小题5分，共40分。在每小题给出的四个选项中，
只有一项是符合题目要求的。}

\begin{enumerate}[itemsep=0.3em]
    \item 题目示例
    \begin{tasks}(4)
        \task 选项A
        \task 选项B
        \task 选项C
        \task 选项D
    \end{tasks}
    \answer{A}
    \analysis{解析内容}
\end{enumerate}

% ══════════════════════════════════════════
%  填空题部分
% ══════════════════════════════════════════

\noindent\textbf{填空题：本题共3小题，每小题5分，共15分。}

\begin{enumerate}[start=9, itemsep=0.8em]
    \item 填空题示例 \blank.
    \answer{答案内容}
    \analysis{解析内容}
\end{enumerate}

% ══════════════════════════════════════════
%  解答题部分
% ══════════════════════════════════════════

\noindent\textbf{解答题：本题共5小题，共77分。解答应写出文字说明、证明过程或演算步骤。}

\begin{examenum}
    \item （13分）解答题示例
    \begin{examenum}
        \item 第1问
        \item 第2问
    \end{examenum}
    \answer{答案内容}
    \analysis{解析内容}
    \explanation{详解内容}
\end{examenum}

\end{document}
